\documentclass[11pt,a4paper]{article}
\usepackage[T1]{fontenc}
\usepackage{lmodern,amsmath,amssymb,booktabs,microtype}
\usepackage[margin=25mm]{geometry}
\usepackage[colorlinks=true,urlcolor=blue,linkcolor=blue]{hyperref}
\newcommand{\tr}{\operatorname{tr}}
\newcommand{\Tr}{\operatorname{Tr}}
\newcommand{\E}{\mathbb E}
\title{Complex DSSYK: a normalized route\\to the first singlet spectral function}
\author{Alan Day \and Chengheng\thanks{AI research assistant (OpenAI).}}
\date{Research note --- 5 November 2026}
\begin{document}\maketitle
\begin{abstract}
We connect the Susskind-line singlet programme to existing solutions of
complex double-scaled SYK and embedded ensembles. A finite-$N$ dictionary
fixes the fermion and coupling normalizations. The first derivative
composite is expressed through neighboring particle-number sectors,
turning its dynamical spectrum into a precisely specified charged
four-point density. Exact ensemble moments and small-system calculations
provide the calibration for the next large-$N$ calculation.
\end{abstract}

\section{Fixing the microscopic model}
Let $N$ count complex modes, $p$ count creation operators per monomial,
and $\lambda=p^2/N$. Work with $\{c_i,c_j^\dagger\}=\delta_{ij}$,
$Q=\sum_i c_i^\dagger c_i$, and the full infinite-temperature trace
$\tr=2^{-N}\Tr$. The fermions in \cite{confine,bnr} have
$\psi_i=\sqrt2 c_i$. Thus the string Hamiltonian of \cite{mss} becomes
\begin{equation}
 H_s=2^p\sum_{|I|=|J|=p}J_{IJ}c_I^\dagger c_J,\qquad
 J_{IJ}=J_{JI}^*,\qquad
 \E J_{IJ}J_{KL}=v_s\delta_{IL}\delta_{JK},
\end{equation}
\begin{equation}
 v_s=\frac{\mathcal J_s^2}{\lambda}
       \left(\frac{p!}{N^p}\right)^2.\label{vs}
\end{equation}
Creation products are ordered and annihilation products are their
adjoints. Reversing both orderings changes no monomial. Diagonal
$J_{II}$ are real with variance $v_s$; off-diagonal entries have
$\E|J_{IJ}|^2=v_s$. In standard-CAR code the coefficient variance is
$4^p v_s$, not $v_s$.

The BNR convention \cite{bnr}, (2.4), uses
$v_B=J_B^2\binom Np^{-2}$. The exact coefficient matching is
\begin{equation}
 J_B=\frac{\mathcal J_s}{\sqrt\lambda}P_{N,p},\qquad
 P_{N,p}=\frac{(N)_p}{N^p}
        =\prod_{j=0}^{p-1}\left(1-\frac jN\right).\label{map}
\end{equation}
At fixed $\lambda>0$, $P_{N,p}\to e^{-\lambda/2}$, since
$\log P_{N,p}=-p(p-1)/(2N)+O(p^3/N^2)$. Keeping this factor is essential
when transferring an existing spectral formula.

\newpage
\section{Reuse the known Hamiltonian solution}
BNR \cite{bnr}, (4.3), gives the zero-chemical-potential chord moments.
After \eqref{map}, their double-scaled form in our units is
\begin{equation}
 \E\tr H_s^{2k}
 =\left(\frac{\mathcal J_s^2}{\lambda}\right)^k
   \sum_{\pi\in\mathcal P_2(2k)}q^{\operatorname{cr}(\pi)},
 \qquad q=e^{-4\lambda},\qquad \E\tr H_s^{2k+1}=0.
\end{equation}
Here the equality is in the double-scaled limit and $\mathcal P_2$
denotes pairings. In particular,
\begin{equation}
 \E\tr H_s^2=\frac{\mathcal J_s^2}{\lambda},\qquad
 \frac{\E\tr H_s^4}{(\E\tr H_s^2)^2}=2+e^{-4\lambda}.
\end{equation}
The existing $q$-Hermite measure can therefore be imported directly;
there is no need to rederive the Hamiltonian density of states.

\paragraph{Time units.}
Use $t_s=p\,t_c$ and $H_c=pH_s$, so that $H_st_s=H_ct_c$.
These relations follow the time-unit convention in \cite{qcd,temp}.
For a fixed probe $O$, $\omega_c=p\omega_s$ and
$S_c(\omega_c)=p^{-1}S_s(\omega_c/p)$. A derivative probe also changes:
$M_{n,c}=p^nM_{n,s}$, hence
\begin{equation}
 A_{2,c}=p^2 A_{2,s},\qquad
 S_{22,c}(\omega_c)=p^3 S_{22,s}(\omega_c/p).
\end{equation}
The coupling symbols in different papers also change definitions.
For the same unit Gaussian coefficients, the 2023 complex Hamiltonian
and the 2025 one obey
$H_{2025}=\mathcal J_{2025}H_{2023}/(p\mathcal J_{2023})$.
We fix \eqref{vs} first and perform time conversion only once.
The Majorana paper \cite{qcd} instead uses total degree $p_M$ and
$\lambda_M=2p_M^2/N_M$; its coefficients are not those of the complex model.

\paragraph{Finite-$N$ normalization.}
Set $w=4^pv_s$. In the $m$-particle sector, of dimension $d_m=\binom Nm$,
the standard embedded-ensemble identity \cite{tt}, (2.8)--(2.9), gives
\begin{equation}
 \E(H_s^2)|_m=F_m I_m,\qquad
 F_m=w\binom mp\binom{N-m+p}{p},\qquad F_m=0\ (m<p).\label{fm}
\end{equation}
The full trace second moment is exactly
$2^{-N}\sum_m d_mF_m$. Equivalently, counting the overlap $k=|I\cap J|$,
\begin{equation}
 \E\tr H_s^2=v_s\binom Np
 \sum_{k=\max(0,2p-N)}^p\binom pk\binom{N-p}{p-k}2^k.\label{overlap}
\end{equation}
Equations \eqref{fm} and \eqref{overlap} are independent implementations
of the same finite-system normalization, used below as a check.

\newpage
\section{The first singlet connects neighboring charge sectors}
Use $O(t)=e^{iH_st}Oe^{-iH_st}$ and
$M_n=\sum_i c_i^\dagger(i\,\operatorname{ad}_{H_s})^n c_i$.
As in the programme note, $M_0=Q$, $M_1=-ipH_s$ and
\begin{equation}
 A_2=-M_2=\sum_i[H_s,c_i]^\dagger[H_s,c_i].
\end{equation}
Define $\mathcal T(X)=\sum_i c_i^\dagger Xc_i$. The homogeneous
normal-ordered CAR identity gives $\mathcal T(H_s)=(Q-p)H_s$.
Expanding the two commutators then yields the useful exact formula
\begin{equation}
 \boxed{A_2=\mathcal T(H_s^2)-(Q-2p)H_s^2.}\label{a2}
\end{equation}
Let $H_m$ denote $H_s$ restricted to particle number $m$. Since $c_i$
lowers particle number by one,
\begin{equation}
 (A_2)_m=\sum_i c_i^\dagger H_{m-1}^2c_i-(m-2p)H_m^2.\label{block}
\end{equation}
The second term is conserved. All nonzero-frequency weight comes from
the first term, which retains the same couplings in adjacent sectors.
For energy eigenstates $r,s$ in sector $m$ and $u$ in sector $m-1$,
\begin{equation}
 (A_2)^{(m)}_{rs}=
 \sum_{i,u}(E_u^{(m-1)})^2(c_i)_{ur}^*(c_i)_{us}
 -(m-2p)(E_r^{(m)})^2\delta_{rs}.\label{matrix}
\end{equation}
This identifies the matrix elements needed for the singlet spectrum.
The coherent sum over $i,u$ must be formed before taking its modulus
squared. The Hamiltonian density of states alone does not supply it.

\paragraph{Two useful exact consequences.}
At the first active sector, $m=p$, the preceding Hamiltonian vanishes,
so $(A_2)_p=pH_p^2$. This sector has only zero-frequency weight.
For $p=1$, $H=c^\dagger h c$ and $A_2=c^\dagger h^2c$ commute in every
sector. These give two distinct conserved controls for a calculation.

Taking the ensemble expectation in \eqref{block} and using \eqref{fm}
also fixes the mean without numerical disorder sampling:
\begin{align}
 \E(A_2)|_m
 &=\big[mF_{m-1}-(m-2p)F_m\big]I_m\nonumber\\
 &=F_m\frac{p(N-p+1)}{N-m+1}I_m,\qquad m\ge p.\label{amean}
\end{align}
The elementary binomial ratio proves the second line. This is a
normalization and static-background formula, not the dynamical spectrum.

\newpage
\section{The precise four-point object to compute next}
Write $D_m(E)=\delta(H_m-E)$, understood as a finite-system spectral
measure. Define the charged four-point density
\begin{align}
 \Xi_m(E,E';u,v)=\frac{1}{d_m}\E\sum_{i,j}\Tr_m\big[&
 c_i^\dagger D_{m-1}(u)c_i D_m(E')\nonumber\\[-2pt]
 &\cdot c_j^\dagger D_{m-1}(v)c_j D_m(E)\big].\label{xi}
\end{align}
In this equation $u,v$ are energy variables, not state labels.
The first probe's nonzero-frequency trace spectrum is exactly
\begin{align}
 S_{22}(\omega)=2\pi\sum_{m=1}^N\frac{d_m}{2^N}
 \int &dE\,dE'\,du\,dv\;u^2v^2\nonumber\\[-2pt]
 &\cdot\Xi_m(E,E';u,v)\delta(\omega-E'+E),\qquad\omega\ne0.\label{target}
\end{align}
Centering $A_2$ changes only its zero-frequency contribution. Likewise,
the conserved term in \eqref{block} drops out of \eqref{target}. The
formula follows by inserting energy projectors in \eqref{matrix};
it is an exact finite-$N$ reformulation of the proposed observable.

\paragraph{Why this is an efficient next target.}
BNR supplies oriented-chord methods for complex fermions, while
Tall--Tomsovic supplies fixed-charge moments and four-point strength
densities \cite{bnr,tt}. We can use those methods to evaluate
\eqref{xi}, retaining both adjacent sectors and the prescribed index
contractions. This is more specific than constructing an unspecified
new ladder kernel from scratch.

The published closed probe formulas in BNR \S7--8 and
Tall--Tomsovic \S4 use independent Gaussian probe coefficients.
Our $A_2$ depends quadratically on the same $J$ as $H_s$.
Equation \eqref{xi} keeps that dependence through explicit elementary
fermion insertions, rather than replacing $A_2$ by an independent
random operator. Individual elementary insertions have fixed length;
their double-scaled limit and the weighted energy integrals must be
controlled together.

\paragraph{The calculation sequence.}
First derive low energy-weighted moments of \eqref{xi} using the
fixed-charge operator identities and oriented contractions. Reproduce
the conserved channels and \eqref{amean}. Then resum the dynamical
channel, starting with small positive $\lambda$ after double scaling.
Finally compare the resulting boost-frequency kernel with Rindler
bulk propagation to extract invariant particle masses and weights.

\newpage
\section{Finite-system calibration}
We implemented \eqref{vs} with $\mathcal J_s=1$ and the exact full
trace. Five covariance sums over all pairs of Hamiltonian monomials
agree with \eqref{overlap} and the mean of $A_2$; an independent
particle-sector sum using \eqref{fm}--\eqref{amean} agrees as well.
These ensemble checks require no Monte Carlo estimate.

Eighteen same-sample calculations use three seeds (1105--1107) for
each pair $(N,p)$ below. Diagonalizing within exact charge sectors,
we store every transition frequency and weight, plus the couplings,
Hamiltonian and $A_2$ matrices. Let $O_2=A_2-\tr(A_2)I$ and
\begin{equation}
 G=\tr O_2^2,\qquad K=\tr([H_s,O_2]^\dagger[H_s,O_2]),\qquad
 f_{\rm dyn}=\frac{\int_{\omega\ne0}S_{22}(\omega)d\omega}{2\pi G}.
\end{equation}
The observed ranges of $f_{\rm dyn}$ are
\begin{center}
\begin{tabular}{rrl}
\toprule $N$&$p$& $f_{\rm dyn}$ across three samples\\\midrule
4&1& below $4\times10^{-28}$ (conserved)\\
3&2& below $3\times10^{-31}$ (conserved)\\
4&2& $0.02414$--$0.07011$\\
5&2& $0.05992$--$0.10605$\\
6&2& $0.12430$--$0.15191$\\
6&3& $0.03598$--$0.03757$\\\bottomrule
\end{tabular}
\end{center}
Thus the construction distinguishes the conserved controls from
interacting samples with nonzero dynamical weight. These small systems
calibrate the observable; they are not an extrapolation of its
double-scaled spectrum. The 178 algebraic and spectral checks include
$M_1$, $M_2$, \eqref{block}, the $m=p$ sector and both spectral sum rules.
The largest recorded scaled residual is $1.30\times10^{-15}$.

\appendix
\section{Reproduction and remaining scope}
The implementation and read-only audit scripts are listed in the
accompanying source ledger and stored in \texttt{code/}.
The environment is Python 3.14.5 and NumPy 2.4.6, also recorded
in \texttt{calibration.json}.
All calculations use complex double precision and save code/data hashes.
The numerical zero bin is $10^{-10}\max(1,\|H_s\|_2)$; reducing it by
ten changes no reported static weight at the test tolerance.

The algebraic formulas are standard CAR/ensemble consequences; the
Hamiltonian chord solution is due to the cited literature. The
large-$N$ evaluation of \eqref{xi}, its spectral normalization and its
bulk mass map remain the next research work. Late-time exponential
approximations to elementary correlators are not differentiated at
zero time to obtain the energy-weighted insertions above.

\newpage
\begin{thebibliography}{9}\small
\bibitem{confine} L. Susskind, \emph{De Sitter Space has no Chords.
Almost Everything is Confined},
\href{https://arxiv.org/abs/2303.00792v1}{arXiv:2303.00792v1}.
\bibitem{mss} S. Miyashita, Y. Sekino and L. Susskind,
\emph{DSSYK at Infinite Temperature: The Flat-Space Limit and the
't Hooft Model}, \href{https://arxiv.org/abs/2506.18054v3}{arXiv:2506.18054v3}.
\bibitem{bnr} M. Berkooz, V. Narovlansky and H. Raj,
\emph{Complex Sachdev-Ye-Kitaev model in the double scaling limit},
JHEP 02 (2021) 113,
\href{https://arxiv.org/abs/2006.13983v2}{arXiv:2006.13983v2}.
\bibitem{qcd} Y. Sekino and L. Susskind, \emph{Double-Scaled SYK,
QCD, and the Flat Space Limit of de Sitter Space},
\href{https://arxiv.org/abs/2501.09423v2}{arXiv:2501.09423v2}.
\bibitem{temp} A. Rahman and L. Susskind, \emph{Infinite Temperature
is Not So Infinite: The Many Temperatures of de Sitter Space},
\href{https://arxiv.org/abs/2401.08555v3}{arXiv:2401.08555v3}.
\bibitem{tt} J. Tall and S. Tomsovic, \emph{Double-scaled bosonic and
fermionic embedded ensembles, complex SYK, and the dual Hilbert space},
\href{https://arxiv.org/abs/2604.14522v2}{arXiv:2604.14522v2}.
\end{thebibliography}
\end{document}
